\begin{figure}[htbp]
\centering
\begin{tikzpicture}[fd,
  t3 hadron/.style={draw=black!70,line width=1.5pt,line cap=butt},
  t3 heavy/.style={draw,solid,line width=1.3pt,fd arrow=#1},
  t3 heavy/.default=0.5,
  t3 lab/.style={font=\small},
  t3 note/.style={font=\footnotesize}]
  % ========================= (a) cut gluon ladder =========================
  \begin{scope}
    \node[hadron blob,ellipse,minimum width=3.0cm,minimum height=0.9cm,inner sep=0pt] (P) at (0,0) {};
    \node[hard blob,ellipse,minimum width=2.6cm,minimum height=0.75cm,inner sep=0pt] (H) at (0,3.95) {};
    \draw[t3 hadron] (-2.9,0) -- (P.west);  \draw[t3 hadron] (P.east) -- (2.9,0);
    \draw[photon] (-2.45,4.6) -- (H.170);  \draw[photon] (H.10) -- (2.45,4.6);
    % t-channel gluons (amplitude left, conjugate right)
    \draw[gluon] (-0.9,0.36) -- (-0.9,0.95);  \draw[gluon] (-0.9,0.95) -- (-0.9,1.7);
    \draw[gluon] (-0.9,1.7) -- (-0.9,2.85);   \draw[gluon] (-0.9,2.85) -- (-0.9,3.62);
    \draw[gluon] (0.9,0.36) -- (0.9,0.95);    \draw[gluon] (0.9,0.95) -- (0.9,1.7);
    \draw[gluon] (0.9,1.7) -- (0.9,2.85);     \draw[gluon] (0.9,2.85) -- (0.9,3.62);
    % s-channel (real) gluons = rungs, crossing the cut
    \foreach \y in {0.95,1.7,2.85} {
      \draw[gluon] (-0.9,\y) -- (0.9,\y);
      \node[vertex] at (-0.9,\y) {};  \node[vertex] at (0.9,\y) {}; }
    \node[fill=white,inner sep=1pt,font=\small] at (0.45,2.27) {$\vdots$};
    \draw[cut] (0,-0.7) -- (0,4.95);
    \node[hard blob,ellipse,minimum width=2.6cm,minimum height=0.75cm,inner sep=0pt] at (0,3.95) {$H$};
    \node[hadron blob,ellipse,minimum width=3.0cm,minimum height=0.9cm,inner sep=0pt] at (0,0) {$p$};
    \node[t3 lab,left=2pt] at (-0.98,0.62) {$k_1$};
    \node[t3 lab,left=2pt] at (-0.98,1.32) {$k_2$};
    \node[t3 lab,left=2pt] at (-0.98,3.24) {$k_n$};
    \node[t3 lab,above left=-1pt] at (-2.45,4.6) {$\gamma^*$};
    \node[t3 lab] at (0,-1.15) {(a)\ 小 $x$ 胶子梯形图};
  \end{scope}
  % ===================== (b) k_T factorization =====================
  \begin{scope}[xshift=6.9cm]
    \node[jet blob,ellipse,minimum width=3.0cm,minimum height=0.9cm,inner sep=0pt] (F) at (0,0) {};
    \coordinate (a) at (-0.9,3.45); \coordinate (ap) at (0.9,3.45);
    \coordinate (b) at (-0.9,2.25); \coordinate (bp) at (0.9,2.25);
    \draw[t3 hadron] (-2.9,0) -- (F.west);  \draw[t3 hadron] (F.east) -- (2.9,0);
    \draw[gluon] (-0.9,0.36) -- (b);  \draw[gluon] (0.9,0.36) -- (bp);
    \draw[t3 heavy=0.3] (a) -- (ap);  \draw[t3 heavy] (ap) -- (bp);
    \draw[t3 heavy=0.3] (bp) -- (b);  \draw[t3 heavy] (b) -- (a);
    \draw[photon] (-2.45,4.25) -- (a);  \draw[photon] (ap) -- (2.45,4.25);
    \foreach \p in {a,ap,b,bp} \node[vertex] at (\p) {};
    \draw[cut] (0,-0.7) -- (0,4.75);
    \node[jet blob,ellipse,minimum width=3.0cm,minimum height=0.9cm,inner sep=0pt] at (0,0) {$\mathcal F(x,\boldsymbol k_{\mathrm T})$};
    \path[mom={$k$},mom sep=0.27cm] (-0.9,0.4) -- (b);
    \path[mom={},mom sep=0.27cm] (bp) -- (0.9,0.4);
    \node[t3 note,anchor=west] at (1.4,1.3) {$k=xP+k_{\mathrm T}$};
    \node[t3 note,anchor=west] at (1.4,0.9) {$k^2\simeq-\boldsymbol k_{\mathrm T}^2$};
    \node[t3 lab,above=1pt] at (-0.45,3.45) {$Q$};
    \node[t3 lab,below=1pt] at (-0.45,2.25) {$\bar Q$};
    \node[t3 note,anchor=east] at (-1.0,2.85) {离壳 $\hat\sigma$};
    \node[t3 lab,above left=-1pt] at (-2.45,4.25) {$\gamma^*$};
    \node[t3 lab] at (0,-1.15) {(b)\ $k_{\mathrm T}$ 因子化};
  \end{scope}
\end{tikzpicture}
\caption{小 $x$ 物理。(a)~割开的胶子梯形图：$t$ 道胶子 $k_1,\dots,k_n$ 自质子 $p$ 向硬过程 $H$ 依次传递，横档（实胶子）穿过割线。在领头对数 $\ln(1/x)$ 近似（BFKL）中，
各级动量分数强有序，$x_1\gg x_2\gg\cdots\gg x_n$，而横动量 $\boldsymbol k_{i\mathrm T}$ 不必有序，图的求和给出 $\sum_n c_n\big(\alpha_s\ln\frac1x\big)^n$；
与之对照，DGLAP 的领头对数 $\ln Q^2$ 来自横动量强有序 $k_{1\mathrm T}\ll k_{2\mathrm T}\ll\cdots\ll Q$ 的区域。
(b)~$k_{\mathrm T}$（高能）因子化：最上一级胶子 $k=xP+k_{\mathrm T}$ 离壳（$k^2\simeq-\boldsymbol k_{\mathrm T}^2$），与离壳硬部分（此处为 $\gamma^*g^*\to Q\bar Q$，粗线为重夸克）相连，
下方蓝色团为非积分胶子分布 $\mathcal F(x,\boldsymbol k_{\mathrm T})$（包含 (a) 中的梯形求和）；截面形如
$\sigma\simeq\int\frac{\dd\xi}{\xi}\int\dd^2\boldsymbol k_{\mathrm T}\,\hat\sigma\big(\tfrac{x}{\xi},\boldsymbol k_{\mathrm T};Q\big)\,\mathcal F(\xi,\boldsymbol k_{\mathrm T})$（示意；归一化约定见正文）。
$k_{\mathrm T}$ 因子化只在 $\ln(1/x)$ 的领头对数精度、对特定过程（如重夸克产生）被确立（[P]）；BFKL 方程本身是微扰论高能极限下的（次）领头对数结果；饱和区的非线性演化属 [C]/[P]。}
\label{fig:bfkl-ladder}
\end{figure}
